% Turning-rule experiments: numbers from build/the-trial/variants/numbers.tex (\lab{...})
\input{../build/the-trial/variants/numbers.tex}
\newcommand\vgraphics{../build/the-trial/variants/}
\newcommand\variationspread[3]{% #1 rule, #2 title, #3 text
  \cleartoverso
  \section{#2}
  #3

  \medskip
  \noindent\begin{minipage}[t]{57mm}\centering
  \includegraphics{\vgraphics#1-plan.pdf}\\[1mm]{\figfont\scriptsize from above}\end{minipage}\hfill
  \begin{minipage}[t]{57mm}\centering
  \includegraphics{\vgraphics#1-shaft.pdf}\\[1mm]{\figfont\scriptsize mineshaft, side view, vertical scale \lab{#1-vscale}}\end{minipage}
  \clearpage
  \thispagestyle{empty}
  \noindent\includegraphics[width=\textwidth]{\vgraphics#1-stand.png}\\[1mm]
  {\figfont\scriptsize Standing letters, eye level, at the busiest junction (\lab{#1-stand-near} sentences within 8 units).}\\[4mm]
  \includegraphics[width=\textwidth]{\vgraphics#1-high.png}\\[1mm]
  {\figfont\scriptsize The same junction from above and behind.}}

\chapter{Other ways to walk}
\label{ch:variations}

The rule used so far is deliberately crude: every feeling, mild or violent,
turns the walker by the same right angle. There are gentler ways to walk a
book, and each shows something different. Here \emph{The Trial} is walked in
three ways, and each walk is seen from above, as a shaft from the side, and on
the ground among its own letters.

\begin{description}
\item[Quarter turns] the rule of the previous chapter. Every sentence is a
straight line, and every change of mood a right angle.
\item[Graded turns] each sentence turns by how strongly it feels. A faintly
cold sentence leans a little to the right; a bitter one turns almost a right
angle. (A sentence's total score is softened into the range from $-1$ to $+1$
the way the VADER sentiment model does it, then multiplied by
90\textdegree.)
\item[Rolling turns] the walker turns not at the end of sentences but at every
word, by the mood of the last \lab{window} words. A plain stretch runs
straight; a passage of feeling bends into a curve as it is read.
\end{description}

\noindent The three walks agree that this translation of \emph{The Trial} leans
slightly warm when every word is scored, and they differ in how much. Counting
sentences, a mild sentence weighs as much as a violent one; weighing them by
strength, or counting words, softens the lean
(\lab{quarter-net_turns}, \lab{graded-net_turns} and
\lab{rolling-net_turns} full turns to the left). All three inherit the problem
of \emph{like}.

\variationspread{quarter}{Quarter turns}{Seen from above, the quarter-turn walk
is a maze of right angles. On the ground, at the place where most sentences
meet, they cross at right angles and their letters pass through one another.}

\variationspread{graded}{Graded turns}{Weighed by strength, the maze becomes a
coastline. Mild sentences bend the path a little, and strong ones fold it.
Each sentence is still a straight line, but neighbouring sentences now meet at
their own angles, so on the ground the sentences fan out like spokes.}

\variationspread{rolling}{Rolling turns}{Walked word by word, the drawing loses
its corners altogether. A plain stretch runs straight, and a passage of
feeling bends into a curve as tight as the feeling is strong. On the ground,
the words of a sentence follow the curve, and a sentence reads along an arc.}
